% ======================================================================
%  AT-MONO108 — Chapter 9: Quantum Mechanics
%  Canonical Monograph (v2.0) · The Actualization Theory:
%  A Reconstruction of Physics from Difference, Actualization and Spectrum
%
%  Status: COMPLETE — PASS-READY (MONO007 referee-ready architecture)
%  Inputs: Chapter01_Difference.tex ... Chapter08_LockLaw.tex,
%          ATQG_CanonicalMonograph.md, ATQG_Mono007RefereeReadiness.md
%  Method: assembly only — canonical results only, no new physics, no speculation
%  Citations: QG216 (amplitude origin), QG218 (Hilbert origin),
%             QG220 (phase origin), QG223 (complete QG),
%             QG243 (gauge dynamics origin), QG244 (Lagrangian origin),
%             QG286 (difference duality), QG301 (duality complete)
%
%  Usage: % ======================================================================
%  AT-MONO108 — Chapter 9: Quantum Mechanics
%  Canonical Monograph (v2.0) · The Actualization Theory:
%  A Reconstruction of Physics from Difference, Actualization and Spectrum
%
%  Status: COMPLETE — PASS-READY (MONO007 referee-ready architecture)
%  Inputs: Chapter01_Difference.tex ... Chapter08_LockLaw.tex,
%          ATQG_CanonicalMonograph.md, ATQG_Mono007RefereeReadiness.md
%  Method: assembly only — canonical results only, no new physics, no speculation
%  Citations: QG216 (amplitude origin), QG218 (Hilbert origin),
%             QG220 (phase origin), QG223 (complete QG),
%             QG243 (gauge dynamics origin), QG244 (Lagrangian origin),
%             QG286 (difference duality), QG301 (duality complete)
%
%  Usage: % ======================================================================
%  AT-MONO108 — Chapter 9: Quantum Mechanics
%  Canonical Monograph (v2.0) · The Actualization Theory:
%  A Reconstruction of Physics from Difference, Actualization and Spectrum
%
%  Status: COMPLETE — PASS-READY (MONO007 referee-ready architecture)
%  Inputs: Chapter01_Difference.tex ... Chapter08_LockLaw.tex,
%          ATQG_CanonicalMonograph.md, ATQG_Mono007RefereeReadiness.md
%  Method: assembly only — canonical results only, no new physics, no speculation
%  Citations: QG216 (amplitude origin), QG218 (Hilbert origin),
%             QG220 (phase origin), QG223 (complete QG),
%             QG243 (gauge dynamics origin), QG244 (Lagrangian origin),
%             QG286 (difference duality), QG301 (duality complete)
%
%  Usage: % ======================================================================
%  AT-MONO108 — Chapter 9: Quantum Mechanics
%  Canonical Monograph (v2.0) · The Actualization Theory:
%  A Reconstruction of Physics from Difference, Actualization and Spectrum
%
%  Status: COMPLETE — PASS-READY (MONO007 referee-ready architecture)
%  Inputs: Chapter01_Difference.tex ... Chapter08_LockLaw.tex,
%          ATQG_CanonicalMonograph.md, ATQG_Mono007RefereeReadiness.md
%  Method: assembly only — canonical results only, no new physics, no speculation
%  Citations: QG216 (amplitude origin), QG218 (Hilbert origin),
%             QG220 (phase origin), QG223 (complete QG),
%             QG243 (gauge dynamics origin), QG244 (Lagrangian origin),
%             QG286 (difference duality), QG301 (duality complete)
%
%  Usage: \input{Chapter09_QuantumMechanics.tex} after the preamble
%         that defines the theorem environments (or include via the master file).
% ======================================================================

% ---- Theorem environments (define in the master preamble if not present) ----
% \newtheorem{definition}{Definition}[chapter]
% \newtheorem{lemma}[definition]{Lemma}
% \newtheorem{theorem}[definition]{Theorem}
% \newtheorem{corollary}[definition]{Corollary}
% \newtheorem{remark}[definition]{Remark}

\chapter{Quantum Mechanics}
\label{c09:chap:qm}

\section{Introduction}
\label{c09:sec:qm-introduction}

The preceding chapters established the primitives $\{\text{Difference},\eta\}$,
Actualization, the D96 spectrum, and the operator basis through which it is read. This
chapter derives the quantum structure of the theory from that foundation: quantum
mechanics is \emph{emergent} --- a read of the D96 spectrum through the Difference duality
--- not a set of independent postulates.

The central results: the quantum state arises from the spectrum structure; the amplitude
magnitude satisfies $|\psi|^2 = \rho$ --- the counting-measure share is the amplitude
magnitude squared \cite{c09:QG216}; the state carries a phase structure derived from the
actualization cycle \cite{c09:QG220}; the complex structure is forced by the two independent
degrees of freedom of the state \cite{c09:QG218}; interference is phase-dependent; and
\emph{measurement is actualization} --- the collapse is a Born-weighted projection identified
with the actualization event \cite{c09:QG223}. The Born rule is derived, not asserted:
$\sum|\psi|^2 = 1$ is exact by construction, as the normalization of the actualization share
\cite{c09:QG216}.

These results are the canonical end-state of the quantum-origin program
\cite{c09:QG216,c09:QG218,c09:QG220,c09:QG223,c09:QG243,c09:QG244}, consistent with the Difference duality
\cite{c09:QG286,c09:QG301}.

\section{Quantum States from Spectrum}
\label{c09:sec:states-from-spectrum}

\begin{theorem}[Quantum states are spectral reads]
\label{c09:thm:states-spectral}
A \emph{quantum state} is a spectral read: a state of the system assigned by the D96
spectrum structure, carrying the degrees of freedom of the underlying actualization
process. The state structure is determined by the spectrum's mode and multiplicity
structure, not by an independent quantum postulate. Quantum mechanics is therefore
emergent in the theory: its states, amplitudes, and structure are derived reads of the
spectrum, not primitive postulates.
\end{theorem}

\begin{proof}
Quantum states are assigned over the spectrum of Chapter~\ref{c06:chap:d96}: the mode and
multiplicity structures determine the state space, and the Hilbert-origin result
\cite{c09:QG218} derives the state structure from the spectral degrees of freedom. Hence the
quantum states are spectral reads --- emergent assignments over the D96 spectrum, not
independent inputs. The amplitude and phase structure are derived in
Sections~\ref{c09:sec:amplitudes} and \ref{c09:sec:phase}, and the measurement rule in
Section~\ref{c09:sec:measurement}; every component of quantum mechanics is thus emergent from
the spectrum, consistent with the canonical architecture \cite{c09:QG223}.
\end{proof}

\section{Difference Duality (\texorpdfstring{$\rho,\psi$}{rho, psi})}
\label{c09:sec:duality}

The Difference duality structures quantum mechanics: the two faces of Difference --- the
scalar $\rho$ and the tensor $\psi$ --- carry, respectively, the amplitude and the
gravitational/tensor content of the state. The duality \cite{c09:QG286} establishes that
$\rho$ and $\psi$ are the trace and traceless faces of the one Difference
(Chapter~\ref{c01:chap:difference}, Theorem~\ref{c01:thm:duality}); the scalar face is the
counting measure whose share is the amplitude magnitude squared
(Section~\ref{c09:sec:amplitudes}), and the tensor face is the Weyl content read against the
reference $\eta$ (Chapter~\ref{c02:chap:eta}, Theorem~\ref{c02:thm:tensor-read}). The duality
is invariant --- it changes meaning, not numbers \cite{c09:QG301}: the quantum structure is
consistent with it, and no number changes under the duality reinterpretation.

\section{Amplitudes and Probabilities}
\label{c09:sec:amplitudes}

\begin{theorem}[The amplitude magnitude is derived]
\label{c09:thm:amplitude-derived}
The \emph{amplitude magnitude} of a state at generation $k$ is the square root of its
counting-measure share: $|\psi_k| = \sqrt{\rho_k}$, where
$\rho_k = \mu^k/S$ is the actualization share, $\mu$ the branching ratio of the
Galton--Watson process, and $S = \sum_{j<K_{\mathrm{gen}}}\mu^j$, with $K_{\mathrm{gen}}$ the
total number of generations. The amplitude magnitude is derived from
the actualization process:
\begin{equation}
\label{c09:eq:amplitude}
|\psi_k|^2 \;=\; \rho_k \;=\; \frac{\mu^k}{S},
\end{equation}
the counting-measure share of the state. The Born rule $\sum_k|\psi_k|^2 = 1$ holds exactly
by construction, as the normalization of the actualization share: probabilities are the
normalized actualization shares.
\end{theorem}

\begin{proof}
The amplitude-origin result \cite{c09:QG216} derives the amplitude magnitude from the Q-event
process: the path multiplicity to generation $k$ is $\mu^k$, and the normalized share
$\rho_k = \mu^k/S$ is the amplitude magnitude squared. The normalization
$\sum_k|\psi_k|^2 = 1$ is exact by construction for any branching ratio $\mu$; the Born
rule follows from the normalization of the share, so probability is the actualization share
--- the measure, not a separately imposed rule.
\end{proof}

Here $\mu^k$ is the combinatorial path multiplicity of the Galton--Watson tree, namely the
number of root-to-generation-$k$ paths, entering only through the normalized weight
$\rho_k = \mu^k/S$. It is not identified with the conserved Q-event count of
Theorem~\ref{c03:thm:count-conservation}.

\section{Phase Structure}
\label{c09:sec:phase}

\begin{theorem}[The phase is derived from actualization]
\label{c09:thm:phase-derived}
The \emph{quantum phase} of a state at depth $k$ is the circulation phase of the
actualization cycle:
\begin{equation}
\label{c09:eq:phase}
\theta_k \;=\; \frac{2\pi k}{N},
\end{equation}
where $N=96$ is the network periodicity and $k$ the branch depth (the actualization tick).
The phase is fixed by causal ordering and the network periodicity of the $N=96$ attractor.
A quantum state carries exactly two independent real degrees of freedom --- the
magnitude $|\psi| = \sqrt{\rho}$ (the node property) and the phase $\theta$ (the link
property) --- and therefore must be represented as a complex state.
\end{theorem}

\begin{proof}
The phase-origin result \cite{c09:QG220} derives the phase from the Q-event process: causal
ordering fixes the position $k$ (the branch depth, the actualization tick), and the network
periodicity of the circulant ring $C_N$, with $N=96$, fixes the phase quantum
$\Delta\theta = 2\pi/N$ by cycle closure --- $N$ ticks advance $2\pi$, giving uniform
circulation. Hence $\theta_k = 2\pi k/N$ is the circulation phase of the actualization
cycle. The Hilbert-origin result \cite{c09:QG218} establishes the two independent real
degrees of freedom: the magnitude $|\psi| = \sqrt{\rho}$ (from the branching counting
measure, a node property) and the phase $\theta$ (from the $U(1)$ link connection, a link
property). A state with magnitude and phase is a complex state; a real-only state space
would give classical addition with no interference. Hence the complex structure is derived
from the two degrees of freedom.
\end{proof}

The lattice phase $\theta_k = 2\pi k/N$ governs \emph{state evolution}. Mixing and CP
phases are \emph{separate} objects: they are derived from spectral ratios via the
chiral-circulation construction, $\delta = \arcsin(r)$ for a spectral ratio $r$, are
continuous functions of those ratios, and are \emph{not} restricted to the state-phase
lattice.

\section{Interference}
\label{c09:sec:interference}

\begin{theorem}[Interference is phase-dependent]
\label{c09:thm:interference}
Interference in the theory is phase-dependent: the two-path amplitude is
\begin{equation}
\label{c09:eq:interference}
P \;=\; \left|e^{i\theta_1} + e^{i\theta_2}\right|^2
\;=\; 2 + 2\cos(\theta_1 - \theta_2),
\end{equation}
with the phase difference producing the interference pattern. Interference is a derived
consequence of the phase structure: the phase-difference pattern of the
actualization-derived phases.
\end{theorem}

\begin{proof}
The Hilbert-origin result \cite{c09:QG218} establishes the interference formula from the
complex state structure: the two-path amplitude is the sum of the two phase-carrying
amplitudes, whose magnitude squared is $2 + 2\cos(\theta_1 - \theta_2)$. The phase
difference, derived from the actualization cycle (Theorem~\ref{c09:thm:phase-derived}),
produces the interference; a real-only state space would give $P = P_1 + P_2$ with no
interference, so the complex structure is necessary. Hence interference is a derived
consequence of the phase structure, not an independent quantum postulate.
\end{proof}

\section{Measurement as Actualization}
\label{c09:sec:measurement}

\begin{theorem}[Measurement is actualization]
\label{c09:thm:measurement-actualization}
\emph{Measurement} is the collapse of the state to a definite outcome: a Born-weighted
projection identified with the actualization event --- a Q-event is a Born-weighted
projection of the state onto a definite outcome, so no separate collapse postulate or
mechanism is required. Consistent with MONO007 \cite{c09:MONO007}, the measurement result
is presented as the accepted canonical identification: collapse equals actualization. The
theory does not add a new mechanism; it identifies the collapse with the process already
established. No stronger claim is made.
\end{theorem}

\begin{proof}
The complete-quantum-gravity audit \cite{c09:QG223} adjudicates the measurement problem: the
collapse is identified with the Q-event actualization --- a Born-weighted projection to a
definite state. Since the Born weights are the actualization shares
(Theorem~\ref{c09:thm:amplitude-derived}), the projection is weighted by the shares and
performed by the actualization event: measurement is actualization, with no separate
collapse mechanism or postulate.
\end{proof}

\section{Consequences}
\label{c09:sec:qm-consequences}

\begin{lemma}[The quantum dynamics is the actualization flow]
\label{c09:lem:qm-dynamics}
The quantum dynamics of the theory is the actualization flow: the interaction dynamics is
the generator action on the spectral modes, and the Lagrangian is the actualization-flow
action.
\end{lemma}

\begin{proof}
The gauge-dynamics origin \cite{c09:QG243} establishes that the interaction dynamics IS the
generator action on the spectral modes: the D96 gauge generators act on the modes, and an
interaction is the generator's action on the mode. The Lagrangian origin \cite{c09:QG244}
derives the Lagrangian density as the actualization-flow action of the D96 generator
fields, $L = -\frac{1}{4}F^a_{\mu\nu}F^{a\mu\nu} + i\bar{\psi}\gamma^\mu D_\mu\psi -
m\bar{\psi}\psi$. Hence the quantum dynamics is the actualization flow, not an imported
dynamics.
\end{proof}

\begin{theorem}[The conservation structure is universal]
\label{c09:thm:qm-conservation}
The quantum conservation structure is the universal conservation principle: the Born norm,
unitarity, and the Noether currents are manifestations of the one conservation rooted in the
identity of Difference.
\end{theorem}

\begin{proof}
The conservation-principle audit \cite{c09:QG267} establishes that the norm (Born rule),
unitarity, trace, count and Bianchi conservation, and the Noether currents are
manifestations of one deeper principle; count conservation is the definitional identity of
the primitive (Chapter~\ref{c03:chap:actualization},
Theorem~\ref{c03:thm:count-conservation}). Hence the quantum conservation structure is the
universal conservation principle, rooted in the identity of Difference.
\end{proof}

Quantum mechanics is therefore consistent with the canonical architecture: states,
amplitudes, phases, interference, measurement, and dynamics are all derived reads of the
D96 spectrum through the Difference duality and the actualization process
\cite{c09:QG286,c09:QG223}.

\section{Summary}
\label{c09:sec:qm-summary}

This chapter has derived the quantum structure of the theory. Quantum mechanics is
emergent in The Actualization Theory: states are spectral reads of the D96 spectrum
(Theorem~\ref{c09:thm:states-spectral}); the amplitude is the counting-measure share,
$|\psi|^2 = \rho$, with the Born rule exact by construction
(Theorem~\ref{c09:thm:amplitude-derived}); the phase is the circulation phase $2\pi k/N$ of
the actualization cycle, and the state is complex (Theorem~\ref{c09:thm:phase-derived});
interference is the phase-difference pattern (Theorem~\ref{c09:thm:interference});
measurement is the actualization event, with no separate collapse postulate
(Theorem~\ref{c09:thm:measurement-actualization}); and the dynamics is the actualization flow
with universal conservation (Lemma~\ref{c09:lem:qm-dynamics},
Theorem~\ref{c09:thm:qm-conservation}). The duality structures QM
(Section~\ref{c09:sec:duality}): the scalar face carries the amplitude, the tensor face the
gravitational content. The quantum pillar is derived, not postulated. The following
chapters derive gravity and spacetime, the standard model, and cosmology from the same
spectral foundation.

\begin{thebibliography}{99}
\bibitem{c09:QG216} AT-QG Phase 216, \emph{Quantum Amplitude Origin}, Docs/Research.
\bibitem{c09:QG218} AT-QG Phase 218, \emph{Hilbert Space Origin}, Docs/Research.
\bibitem{c09:QG220} AT-QG Phase 220, \emph{Quantum Phase Origin}, Docs/Research.
\bibitem{c09:QG223} AT-QG Phase 223, \emph{Final Quantum Gravity Audit} (Complete QG), Docs/Research.
\bibitem{c09:QG243} AT-QG Phase 243, \emph{Gauge Dynamics Origin}, Docs/Research.
\bibitem{c09:QG244} AT-QG Phase 244, \emph{Lagrangian Origin}, Docs/Research.
\bibitem{c09:QG267} AT-QG Phase 267, \emph{Conservation Principle Audit}, Docs/Research.
\bibitem{c09:QG286} AT-QG Phase 286, \emph{Difference Duality Audit}, Docs/Research.
\bibitem{c09:QG301} AT-QG Phase 301, \emph{Duality Prediction Audit}, Docs/Research.
\bibitem{c09:MONO007} AT-MONO007, \emph{Referee-Readiness Resolution}, Docs/Zenodo/Monograph.
\end{thebibliography}
 after the preamble
%         that defines the theorem environments (or include via the master file).
% ======================================================================

% ---- Theorem environments (define in the master preamble if not present) ----
% \newtheorem{definition}{Definition}[chapter]
% \newtheorem{lemma}[definition]{Lemma}
% \newtheorem{theorem}[definition]{Theorem}
% \newtheorem{corollary}[definition]{Corollary}
% \newtheorem{remark}[definition]{Remark}

\chapter{Quantum Mechanics}
\label{c09:chap:qm}

\section{Introduction}
\label{c09:sec:qm-introduction}

The preceding chapters established the primitives $\{\text{Difference},\eta\}$,
Actualization, the D96 spectrum, and the operator basis through which it is read. This
chapter derives the quantum structure of the theory from that foundation: quantum
mechanics is \emph{emergent} --- a read of the D96 spectrum through the Difference duality
--- not a set of independent postulates.

The central results: the quantum state arises from the spectrum structure; the amplitude
magnitude satisfies $|\psi|^2 = \rho$ --- the counting-measure share is the amplitude
magnitude squared \cite{c09:QG216}; the state carries a phase structure derived from the
actualization cycle \cite{c09:QG220}; the complex structure is forced by the two independent
degrees of freedom of the state \cite{c09:QG218}; interference is phase-dependent; and
\emph{measurement is actualization} --- the collapse is a Born-weighted projection identified
with the actualization event \cite{c09:QG223}. The Born rule is derived, not asserted:
$\sum|\psi|^2 = 1$ is exact by construction, as the normalization of the actualization share
\cite{c09:QG216}.

These results are the canonical end-state of the quantum-origin program
\cite{c09:QG216,c09:QG218,c09:QG220,c09:QG223,c09:QG243,c09:QG244}, consistent with the Difference duality
\cite{c09:QG286,c09:QG301}.

\section{Quantum States from Spectrum}
\label{c09:sec:states-from-spectrum}

\begin{theorem}[Quantum states are spectral reads]
\label{c09:thm:states-spectral}
A \emph{quantum state} is a spectral read: a state of the system assigned by the D96
spectrum structure, carrying the degrees of freedom of the underlying actualization
process. The state structure is determined by the spectrum's mode and multiplicity
structure, not by an independent quantum postulate. Quantum mechanics is therefore
emergent in the theory: its states, amplitudes, and structure are derived reads of the
spectrum, not primitive postulates.
\end{theorem}

\begin{proof}
Quantum states are assigned over the spectrum of Chapter~\ref{c06:chap:d96}: the mode and
multiplicity structures determine the state space, and the Hilbert-origin result
\cite{c09:QG218} derives the state structure from the spectral degrees of freedom. Hence the
quantum states are spectral reads --- emergent assignments over the D96 spectrum, not
independent inputs. The amplitude and phase structure are derived in
Sections~\ref{c09:sec:amplitudes} and \ref{c09:sec:phase}, and the measurement rule in
Section~\ref{c09:sec:measurement}; every component of quantum mechanics is thus emergent from
the spectrum, consistent with the canonical architecture \cite{c09:QG223}.
\end{proof}

\section{Difference Duality (\texorpdfstring{$\rho,\psi$}{rho, psi})}
\label{c09:sec:duality}

The Difference duality structures quantum mechanics: the two faces of Difference --- the
scalar $\rho$ and the tensor $\psi$ --- carry, respectively, the amplitude and the
gravitational/tensor content of the state. The duality \cite{c09:QG286} establishes that
$\rho$ and $\psi$ are the trace and traceless faces of the one Difference
(Chapter~\ref{c01:chap:difference}, Theorem~\ref{c01:thm:duality}); the scalar face is the
counting measure whose share is the amplitude magnitude squared
(Section~\ref{c09:sec:amplitudes}), and the tensor face is the Weyl content read against the
reference $\eta$ (Chapter~\ref{c02:chap:eta}, Theorem~\ref{c02:thm:tensor-read}). The duality
is invariant --- it changes meaning, not numbers \cite{c09:QG301}: the quantum structure is
consistent with it, and no number changes under the duality reinterpretation.

\section{Amplitudes and Probabilities}
\label{c09:sec:amplitudes}

\begin{theorem}[The amplitude magnitude is derived]
\label{c09:thm:amplitude-derived}
The \emph{amplitude magnitude} of a state at generation $k$ is the square root of its
counting-measure share: $|\psi_k| = \sqrt{\rho_k}$, where
$\rho_k = \mu^k/S$ is the actualization share, $\mu$ the branching ratio of the
Galton--Watson process, and $S = \sum_{j<K_{\mathrm{gen}}}\mu^j$, with $K_{\mathrm{gen}}$ the
total number of generations. The amplitude magnitude is derived from
the actualization process:
\begin{equation}
\label{c09:eq:amplitude}
|\psi_k|^2 \;=\; \rho_k \;=\; \frac{\mu^k}{S},
\end{equation}
the counting-measure share of the state. The Born rule $\sum_k|\psi_k|^2 = 1$ holds exactly
by construction, as the normalization of the actualization share: probabilities are the
normalized actualization shares.
\end{theorem}

\begin{proof}
The amplitude-origin result \cite{c09:QG216} derives the amplitude magnitude from the Q-event
process: the path multiplicity to generation $k$ is $\mu^k$, and the normalized share
$\rho_k = \mu^k/S$ is the amplitude magnitude squared. The normalization
$\sum_k|\psi_k|^2 = 1$ is exact by construction for any branching ratio $\mu$; the Born
rule follows from the normalization of the share, so probability is the actualization share
--- the measure, not a separately imposed rule.
\end{proof}

Here $\mu^k$ is the combinatorial path multiplicity of the Galton--Watson tree, namely the
number of root-to-generation-$k$ paths, entering only through the normalized weight
$\rho_k = \mu^k/S$. It is not identified with the conserved Q-event count of
Theorem~\ref{c03:thm:count-conservation}.

\section{Phase Structure}
\label{c09:sec:phase}

\begin{theorem}[The phase is derived from actualization]
\label{c09:thm:phase-derived}
The \emph{quantum phase} of a state at depth $k$ is the circulation phase of the
actualization cycle:
\begin{equation}
\label{c09:eq:phase}
\theta_k \;=\; \frac{2\pi k}{N},
\end{equation}
where $N=96$ is the network periodicity and $k$ the branch depth (the actualization tick).
The phase is fixed by causal ordering and the network periodicity of the $N=96$ attractor.
A quantum state carries exactly two independent real degrees of freedom --- the
magnitude $|\psi| = \sqrt{\rho}$ (the node property) and the phase $\theta$ (the link
property) --- and therefore must be represented as a complex state.
\end{theorem}

\begin{proof}
The phase-origin result \cite{c09:QG220} derives the phase from the Q-event process: causal
ordering fixes the position $k$ (the branch depth, the actualization tick), and the network
periodicity of the circulant ring $C_N$, with $N=96$, fixes the phase quantum
$\Delta\theta = 2\pi/N$ by cycle closure --- $N$ ticks advance $2\pi$, giving uniform
circulation. Hence $\theta_k = 2\pi k/N$ is the circulation phase of the actualization
cycle. The Hilbert-origin result \cite{c09:QG218} establishes the two independent real
degrees of freedom: the magnitude $|\psi| = \sqrt{\rho}$ (from the branching counting
measure, a node property) and the phase $\theta$ (from the $U(1)$ link connection, a link
property). A state with magnitude and phase is a complex state; a real-only state space
would give classical addition with no interference. Hence the complex structure is derived
from the two degrees of freedom.
\end{proof}

The lattice phase $\theta_k = 2\pi k/N$ governs \emph{state evolution}. Mixing and CP
phases are \emph{separate} objects: they are derived from spectral ratios via the
chiral-circulation construction, $\delta = \arcsin(r)$ for a spectral ratio $r$, are
continuous functions of those ratios, and are \emph{not} restricted to the state-phase
lattice.

\section{Interference}
\label{c09:sec:interference}

\begin{theorem}[Interference is phase-dependent]
\label{c09:thm:interference}
Interference in the theory is phase-dependent: the two-path amplitude is
\begin{equation}
\label{c09:eq:interference}
P \;=\; \left|e^{i\theta_1} + e^{i\theta_2}\right|^2
\;=\; 2 + 2\cos(\theta_1 - \theta_2),
\end{equation}
with the phase difference producing the interference pattern. Interference is a derived
consequence of the phase structure: the phase-difference pattern of the
actualization-derived phases.
\end{theorem}

\begin{proof}
The Hilbert-origin result \cite{c09:QG218} establishes the interference formula from the
complex state structure: the two-path amplitude is the sum of the two phase-carrying
amplitudes, whose magnitude squared is $2 + 2\cos(\theta_1 - \theta_2)$. The phase
difference, derived from the actualization cycle (Theorem~\ref{c09:thm:phase-derived}),
produces the interference; a real-only state space would give $P = P_1 + P_2$ with no
interference, so the complex structure is necessary. Hence interference is a derived
consequence of the phase structure, not an independent quantum postulate.
\end{proof}

\section{Measurement as Actualization}
\label{c09:sec:measurement}

\begin{theorem}[Measurement is actualization]
\label{c09:thm:measurement-actualization}
\emph{Measurement} is the collapse of the state to a definite outcome: a Born-weighted
projection identified with the actualization event --- a Q-event is a Born-weighted
projection of the state onto a definite outcome, so no separate collapse postulate or
mechanism is required. Consistent with MONO007 \cite{c09:MONO007}, the measurement result
is presented as the accepted canonical identification: collapse equals actualization. The
theory does not add a new mechanism; it identifies the collapse with the process already
established. No stronger claim is made.
\end{theorem}

\begin{proof}
The complete-quantum-gravity audit \cite{c09:QG223} adjudicates the measurement problem: the
collapse is identified with the Q-event actualization --- a Born-weighted projection to a
definite state. Since the Born weights are the actualization shares
(Theorem~\ref{c09:thm:amplitude-derived}), the projection is weighted by the shares and
performed by the actualization event: measurement is actualization, with no separate
collapse mechanism or postulate.
\end{proof}

\section{Consequences}
\label{c09:sec:qm-consequences}

\begin{lemma}[The quantum dynamics is the actualization flow]
\label{c09:lem:qm-dynamics}
The quantum dynamics of the theory is the actualization flow: the interaction dynamics is
the generator action on the spectral modes, and the Lagrangian is the actualization-flow
action.
\end{lemma}

\begin{proof}
The gauge-dynamics origin \cite{c09:QG243} establishes that the interaction dynamics IS the
generator action on the spectral modes: the D96 gauge generators act on the modes, and an
interaction is the generator's action on the mode. The Lagrangian origin \cite{c09:QG244}
derives the Lagrangian density as the actualization-flow action of the D96 generator
fields, $L = -\frac{1}{4}F^a_{\mu\nu}F^{a\mu\nu} + i\bar{\psi}\gamma^\mu D_\mu\psi -
m\bar{\psi}\psi$. Hence the quantum dynamics is the actualization flow, not an imported
dynamics.
\end{proof}

\begin{theorem}[The conservation structure is universal]
\label{c09:thm:qm-conservation}
The quantum conservation structure is the universal conservation principle: the Born norm,
unitarity, and the Noether currents are manifestations of the one conservation rooted in the
identity of Difference.
\end{theorem}

\begin{proof}
The conservation-principle audit \cite{c09:QG267} establishes that the norm (Born rule),
unitarity, trace, count and Bianchi conservation, and the Noether currents are
manifestations of one deeper principle; count conservation is the definitional identity of
the primitive (Chapter~\ref{c03:chap:actualization},
Theorem~\ref{c03:thm:count-conservation}). Hence the quantum conservation structure is the
universal conservation principle, rooted in the identity of Difference.
\end{proof}

Quantum mechanics is therefore consistent with the canonical architecture: states,
amplitudes, phases, interference, measurement, and dynamics are all derived reads of the
D96 spectrum through the Difference duality and the actualization process
\cite{c09:QG286,c09:QG223}.

\section{Summary}
\label{c09:sec:qm-summary}

This chapter has derived the quantum structure of the theory. Quantum mechanics is
emergent in The Actualization Theory: states are spectral reads of the D96 spectrum
(Theorem~\ref{c09:thm:states-spectral}); the amplitude is the counting-measure share,
$|\psi|^2 = \rho$, with the Born rule exact by construction
(Theorem~\ref{c09:thm:amplitude-derived}); the phase is the circulation phase $2\pi k/N$ of
the actualization cycle, and the state is complex (Theorem~\ref{c09:thm:phase-derived});
interference is the phase-difference pattern (Theorem~\ref{c09:thm:interference});
measurement is the actualization event, with no separate collapse postulate
(Theorem~\ref{c09:thm:measurement-actualization}); and the dynamics is the actualization flow
with universal conservation (Lemma~\ref{c09:lem:qm-dynamics},
Theorem~\ref{c09:thm:qm-conservation}). The duality structures QM
(Section~\ref{c09:sec:duality}): the scalar face carries the amplitude, the tensor face the
gravitational content. The quantum pillar is derived, not postulated. The following
chapters derive gravity and spacetime, the standard model, and cosmology from the same
spectral foundation.

\begin{thebibliography}{99}
\bibitem{c09:QG216} AT-QG Phase 216, \emph{Quantum Amplitude Origin}, Docs/Research.
\bibitem{c09:QG218} AT-QG Phase 218, \emph{Hilbert Space Origin}, Docs/Research.
\bibitem{c09:QG220} AT-QG Phase 220, \emph{Quantum Phase Origin}, Docs/Research.
\bibitem{c09:QG223} AT-QG Phase 223, \emph{Final Quantum Gravity Audit} (Complete QG), Docs/Research.
\bibitem{c09:QG243} AT-QG Phase 243, \emph{Gauge Dynamics Origin}, Docs/Research.
\bibitem{c09:QG244} AT-QG Phase 244, \emph{Lagrangian Origin}, Docs/Research.
\bibitem{c09:QG267} AT-QG Phase 267, \emph{Conservation Principle Audit}, Docs/Research.
\bibitem{c09:QG286} AT-QG Phase 286, \emph{Difference Duality Audit}, Docs/Research.
\bibitem{c09:QG301} AT-QG Phase 301, \emph{Duality Prediction Audit}, Docs/Research.
\bibitem{c09:MONO007} AT-MONO007, \emph{Referee-Readiness Resolution}, Docs/Zenodo/Monograph.
\end{thebibliography}
 after the preamble
%         that defines the theorem environments (or include via the master file).
% ======================================================================

% ---- Theorem environments (define in the master preamble if not present) ----
% \newtheorem{definition}{Definition}[chapter]
% \newtheorem{lemma}[definition]{Lemma}
% \newtheorem{theorem}[definition]{Theorem}
% \newtheorem{corollary}[definition]{Corollary}
% \newtheorem{remark}[definition]{Remark}

\chapter{Quantum Mechanics}
\label{c09:chap:qm}

\section{Introduction}
\label{c09:sec:qm-introduction}

The preceding chapters established the primitives $\{\text{Difference},\eta\}$,
Actualization, the D96 spectrum, and the operator basis through which it is read. This
chapter derives the quantum structure of the theory from that foundation: quantum
mechanics is \emph{emergent} --- a read of the D96 spectrum through the Difference duality
--- not a set of independent postulates.

The central results: the quantum state arises from the spectrum structure; the amplitude
magnitude satisfies $|\psi|^2 = \rho$ --- the counting-measure share is the amplitude
magnitude squared \cite{c09:QG216}; the state carries a phase structure derived from the
actualization cycle \cite{c09:QG220}; the complex structure is forced by the two independent
degrees of freedom of the state \cite{c09:QG218}; interference is phase-dependent; and
\emph{measurement is actualization} --- the collapse is a Born-weighted projection identified
with the actualization event \cite{c09:QG223}. The Born rule is derived, not asserted:
$\sum|\psi|^2 = 1$ is exact by construction, as the normalization of the actualization share
\cite{c09:QG216}.

These results are the canonical end-state of the quantum-origin program
\cite{c09:QG216,c09:QG218,c09:QG220,c09:QG223,c09:QG243,c09:QG244}, consistent with the Difference duality
\cite{c09:QG286,c09:QG301}.

\section{Quantum States from Spectrum}
\label{c09:sec:states-from-spectrum}

\begin{theorem}[Quantum states are spectral reads]
\label{c09:thm:states-spectral}
A \emph{quantum state} is a spectral read: a state of the system assigned by the D96
spectrum structure, carrying the degrees of freedom of the underlying actualization
process. The state structure is determined by the spectrum's mode and multiplicity
structure, not by an independent quantum postulate. Quantum mechanics is therefore
emergent in the theory: its states, amplitudes, and structure are derived reads of the
spectrum, not primitive postulates.
\end{theorem}

\begin{proof}
Quantum states are assigned over the spectrum of Chapter~\ref{c06:chap:d96}: the mode and
multiplicity structures determine the state space, and the Hilbert-origin result
\cite{c09:QG218} derives the state structure from the spectral degrees of freedom. Hence the
quantum states are spectral reads --- emergent assignments over the D96 spectrum, not
independent inputs. The amplitude and phase structure are derived in
Sections~\ref{c09:sec:amplitudes} and \ref{c09:sec:phase}, and the measurement rule in
Section~\ref{c09:sec:measurement}; every component of quantum mechanics is thus emergent from
the spectrum, consistent with the canonical architecture \cite{c09:QG223}.
\end{proof}

\section{Difference Duality (\texorpdfstring{$\rho,\psi$}{rho, psi})}
\label{c09:sec:duality}

The Difference duality structures quantum mechanics: the two faces of Difference --- the
scalar $\rho$ and the tensor $\psi$ --- carry, respectively, the amplitude and the
gravitational/tensor content of the state. The duality \cite{c09:QG286} establishes that
$\rho$ and $\psi$ are the trace and traceless faces of the one Difference
(Chapter~\ref{c01:chap:difference}, Theorem~\ref{c01:thm:duality}); the scalar face is the
counting measure whose share is the amplitude magnitude squared
(Section~\ref{c09:sec:amplitudes}), and the tensor face is the Weyl content read against the
reference $\eta$ (Chapter~\ref{c02:chap:eta}, Theorem~\ref{c02:thm:tensor-read}). The duality
is invariant --- it changes meaning, not numbers \cite{c09:QG301}: the quantum structure is
consistent with it, and no number changes under the duality reinterpretation.

\section{Amplitudes and Probabilities}
\label{c09:sec:amplitudes}

\begin{theorem}[The amplitude magnitude is derived]
\label{c09:thm:amplitude-derived}
The \emph{amplitude magnitude} of a state at generation $k$ is the square root of its
counting-measure share: $|\psi_k| = \sqrt{\rho_k}$, where
$\rho_k = \mu^k/S$ is the actualization share, $\mu$ the branching ratio of the
Galton--Watson process, and $S = \sum_{j<K_{\mathrm{gen}}}\mu^j$, with $K_{\mathrm{gen}}$ the
total number of generations. The amplitude magnitude is derived from
the actualization process:
\begin{equation}
\label{c09:eq:amplitude}
|\psi_k|^2 \;=\; \rho_k \;=\; \frac{\mu^k}{S},
\end{equation}
the counting-measure share of the state. The Born rule $\sum_k|\psi_k|^2 = 1$ holds exactly
by construction, as the normalization of the actualization share: probabilities are the
normalized actualization shares.
\end{theorem}

\begin{proof}
The amplitude-origin result \cite{c09:QG216} derives the amplitude magnitude from the Q-event
process: the path multiplicity to generation $k$ is $\mu^k$, and the normalized share
$\rho_k = \mu^k/S$ is the amplitude magnitude squared. The normalization
$\sum_k|\psi_k|^2 = 1$ is exact by construction for any branching ratio $\mu$; the Born
rule follows from the normalization of the share, so probability is the actualization share
--- the measure, not a separately imposed rule.
\end{proof}

Here $\mu^k$ is the combinatorial path multiplicity of the Galton--Watson tree, namely the
number of root-to-generation-$k$ paths, entering only through the normalized weight
$\rho_k = \mu^k/S$. It is not identified with the conserved Q-event count of
Theorem~\ref{c03:thm:count-conservation}.

\section{Phase Structure}
\label{c09:sec:phase}

\begin{theorem}[The phase is derived from actualization]
\label{c09:thm:phase-derived}
The \emph{quantum phase} of a state at depth $k$ is the circulation phase of the
actualization cycle:
\begin{equation}
\label{c09:eq:phase}
\theta_k \;=\; \frac{2\pi k}{N},
\end{equation}
where $N=96$ is the network periodicity and $k$ the branch depth (the actualization tick).
The phase is fixed by causal ordering and the network periodicity of the $N=96$ attractor.
A quantum state carries exactly two independent real degrees of freedom --- the
magnitude $|\psi| = \sqrt{\rho}$ (the node property) and the phase $\theta$ (the link
property) --- and therefore must be represented as a complex state.
\end{theorem}

\begin{proof}
The phase-origin result \cite{c09:QG220} derives the phase from the Q-event process: causal
ordering fixes the position $k$ (the branch depth, the actualization tick), and the network
periodicity of the circulant ring $C_N$, with $N=96$, fixes the phase quantum
$\Delta\theta = 2\pi/N$ by cycle closure --- $N$ ticks advance $2\pi$, giving uniform
circulation. Hence $\theta_k = 2\pi k/N$ is the circulation phase of the actualization
cycle. The Hilbert-origin result \cite{c09:QG218} establishes the two independent real
degrees of freedom: the magnitude $|\psi| = \sqrt{\rho}$ (from the branching counting
measure, a node property) and the phase $\theta$ (from the $U(1)$ link connection, a link
property). A state with magnitude and phase is a complex state; a real-only state space
would give classical addition with no interference. Hence the complex structure is derived
from the two degrees of freedom.
\end{proof}

The lattice phase $\theta_k = 2\pi k/N$ governs \emph{state evolution}. Mixing and CP
phases are \emph{separate} objects: they are derived from spectral ratios via the
chiral-circulation construction, $\delta = \arcsin(r)$ for a spectral ratio $r$, are
continuous functions of those ratios, and are \emph{not} restricted to the state-phase
lattice.

\section{Interference}
\label{c09:sec:interference}

\begin{theorem}[Interference is phase-dependent]
\label{c09:thm:interference}
Interference in the theory is phase-dependent: the two-path amplitude is
\begin{equation}
\label{c09:eq:interference}
P \;=\; \left|e^{i\theta_1} + e^{i\theta_2}\right|^2
\;=\; 2 + 2\cos(\theta_1 - \theta_2),
\end{equation}
with the phase difference producing the interference pattern. Interference is a derived
consequence of the phase structure: the phase-difference pattern of the
actualization-derived phases.
\end{theorem}

\begin{proof}
The Hilbert-origin result \cite{c09:QG218} establishes the interference formula from the
complex state structure: the two-path amplitude is the sum of the two phase-carrying
amplitudes, whose magnitude squared is $2 + 2\cos(\theta_1 - \theta_2)$. The phase
difference, derived from the actualization cycle (Theorem~\ref{c09:thm:phase-derived}),
produces the interference; a real-only state space would give $P = P_1 + P_2$ with no
interference, so the complex structure is necessary. Hence interference is a derived
consequence of the phase structure, not an independent quantum postulate.
\end{proof}

\section{Measurement as Actualization}
\label{c09:sec:measurement}

\begin{theorem}[Measurement is actualization]
\label{c09:thm:measurement-actualization}
\emph{Measurement} is the collapse of the state to a definite outcome: a Born-weighted
projection identified with the actualization event --- a Q-event is a Born-weighted
projection of the state onto a definite outcome, so no separate collapse postulate or
mechanism is required. Consistent with MONO007 \cite{c09:MONO007}, the measurement result
is presented as the accepted canonical identification: collapse equals actualization. The
theory does not add a new mechanism; it identifies the collapse with the process already
established. No stronger claim is made.
\end{theorem}

\begin{proof}
The complete-quantum-gravity audit \cite{c09:QG223} adjudicates the measurement problem: the
collapse is identified with the Q-event actualization --- a Born-weighted projection to a
definite state. Since the Born weights are the actualization shares
(Theorem~\ref{c09:thm:amplitude-derived}), the projection is weighted by the shares and
performed by the actualization event: measurement is actualization, with no separate
collapse mechanism or postulate.
\end{proof}

\section{Consequences}
\label{c09:sec:qm-consequences}

\begin{lemma}[The quantum dynamics is the actualization flow]
\label{c09:lem:qm-dynamics}
The quantum dynamics of the theory is the actualization flow: the interaction dynamics is
the generator action on the spectral modes, and the Lagrangian is the actualization-flow
action.
\end{lemma}

\begin{proof}
The gauge-dynamics origin \cite{c09:QG243} establishes that the interaction dynamics IS the
generator action on the spectral modes: the D96 gauge generators act on the modes, and an
interaction is the generator's action on the mode. The Lagrangian origin \cite{c09:QG244}
derives the Lagrangian density as the actualization-flow action of the D96 generator
fields, $L = -\frac{1}{4}F^a_{\mu\nu}F^{a\mu\nu} + i\bar{\psi}\gamma^\mu D_\mu\psi -
m\bar{\psi}\psi$. Hence the quantum dynamics is the actualization flow, not an imported
dynamics.
\end{proof}

\begin{theorem}[The conservation structure is universal]
\label{c09:thm:qm-conservation}
The quantum conservation structure is the universal conservation principle: the Born norm,
unitarity, and the Noether currents are manifestations of the one conservation rooted in the
identity of Difference.
\end{theorem}

\begin{proof}
The conservation-principle audit \cite{c09:QG267} establishes that the norm (Born rule),
unitarity, trace, count and Bianchi conservation, and the Noether currents are
manifestations of one deeper principle; count conservation is the definitional identity of
the primitive (Chapter~\ref{c03:chap:actualization},
Theorem~\ref{c03:thm:count-conservation}). Hence the quantum conservation structure is the
universal conservation principle, rooted in the identity of Difference.
\end{proof}

Quantum mechanics is therefore consistent with the canonical architecture: states,
amplitudes, phases, interference, measurement, and dynamics are all derived reads of the
D96 spectrum through the Difference duality and the actualization process
\cite{c09:QG286,c09:QG223}.

\section{Summary}
\label{c09:sec:qm-summary}

This chapter has derived the quantum structure of the theory. Quantum mechanics is
emergent in The Actualization Theory: states are spectral reads of the D96 spectrum
(Theorem~\ref{c09:thm:states-spectral}); the amplitude is the counting-measure share,
$|\psi|^2 = \rho$, with the Born rule exact by construction
(Theorem~\ref{c09:thm:amplitude-derived}); the phase is the circulation phase $2\pi k/N$ of
the actualization cycle, and the state is complex (Theorem~\ref{c09:thm:phase-derived});
interference is the phase-difference pattern (Theorem~\ref{c09:thm:interference});
measurement is the actualization event, with no separate collapse postulate
(Theorem~\ref{c09:thm:measurement-actualization}); and the dynamics is the actualization flow
with universal conservation (Lemma~\ref{c09:lem:qm-dynamics},
Theorem~\ref{c09:thm:qm-conservation}). The duality structures QM
(Section~\ref{c09:sec:duality}): the scalar face carries the amplitude, the tensor face the
gravitational content. The quantum pillar is derived, not postulated. The following
chapters derive gravity and spacetime, the standard model, and cosmology from the same
spectral foundation.

\begin{thebibliography}{99}
\bibitem{c09:QG216} AT-QG Phase 216, \emph{Quantum Amplitude Origin}, Docs/Research.
\bibitem{c09:QG218} AT-QG Phase 218, \emph{Hilbert Space Origin}, Docs/Research.
\bibitem{c09:QG220} AT-QG Phase 220, \emph{Quantum Phase Origin}, Docs/Research.
\bibitem{c09:QG223} AT-QG Phase 223, \emph{Final Quantum Gravity Audit} (Complete QG), Docs/Research.
\bibitem{c09:QG243} AT-QG Phase 243, \emph{Gauge Dynamics Origin}, Docs/Research.
\bibitem{c09:QG244} AT-QG Phase 244, \emph{Lagrangian Origin}, Docs/Research.
\bibitem{c09:QG267} AT-QG Phase 267, \emph{Conservation Principle Audit}, Docs/Research.
\bibitem{c09:QG286} AT-QG Phase 286, \emph{Difference Duality Audit}, Docs/Research.
\bibitem{c09:QG301} AT-QG Phase 301, \emph{Duality Prediction Audit}, Docs/Research.
\bibitem{c09:MONO007} AT-MONO007, \emph{Referee-Readiness Resolution}, Docs/Zenodo/Monograph.
\end{thebibliography}
 after the preamble
%         that defines the theorem environments (or include via the master file).
% ======================================================================

% ---- Theorem environments (define in the master preamble if not present) ----
% \newtheorem{definition}{Definition}[chapter]
% \newtheorem{lemma}[definition]{Lemma}
% \newtheorem{theorem}[definition]{Theorem}
% \newtheorem{corollary}[definition]{Corollary}
% \newtheorem{remark}[definition]{Remark}

\chapter{Quantum Mechanics}
\label{c09:chap:qm}

\section{Introduction}
\label{c09:sec:qm-introduction}

The preceding chapters established the primitives $\{\text{Difference},\eta\}$,
Actualization, the D96 spectrum, and the operator basis through which it is read. This
chapter derives the quantum structure of the theory from that foundation: quantum
mechanics is \emph{emergent} --- a read of the D96 spectrum through the Difference duality
--- not a set of independent postulates.

The central results: the quantum state arises from the spectrum structure; the amplitude
magnitude satisfies $|\psi|^2 = \rho$ --- the counting-measure share is the amplitude
magnitude squared \cite{c09:QG216}; the state carries a phase structure derived from the
actualization cycle \cite{c09:QG220}; the complex structure is forced by the two independent
degrees of freedom of the state \cite{c09:QG218}; interference is phase-dependent; and
\emph{measurement is actualization} --- the collapse is a Born-weighted projection identified
with the actualization event \cite{c09:QG223}. The Born rule is derived, not asserted:
$\sum|\psi|^2 = 1$ is exact by construction, as the normalization of the actualization share
\cite{c09:QG216}.

These results are the canonical end-state of the quantum-origin program
\cite{c09:QG216,c09:QG218,c09:QG220,c09:QG223,c09:QG243,c09:QG244}, consistent with the Difference duality
\cite{c09:QG286,c09:QG301}.

\section{Quantum States from Spectrum}
\label{c09:sec:states-from-spectrum}

\begin{theorem}[Quantum states are spectral reads]
\label{c09:thm:states-spectral}
A \emph{quantum state} is a spectral read: a state of the system assigned by the D96
spectrum structure, carrying the degrees of freedom of the underlying actualization
process. The state structure is determined by the spectrum's mode and multiplicity
structure, not by an independent quantum postulate. Quantum mechanics is therefore
emergent in the theory: its states, amplitudes, and structure are derived reads of the
spectrum, not primitive postulates.
\end{theorem}

\begin{proof}
Quantum states are assigned over the spectrum of Chapter~\ref{c06:chap:d96}: the mode and
multiplicity structures determine the state space, and the Hilbert-origin result
\cite{c09:QG218} derives the state structure from the spectral degrees of freedom. Hence the
quantum states are spectral reads --- emergent assignments over the D96 spectrum, not
independent inputs. The amplitude and phase structure are derived in
Sections~\ref{c09:sec:amplitudes} and \ref{c09:sec:phase}, and the measurement rule in
Section~\ref{c09:sec:measurement}; every component of quantum mechanics is thus emergent from
the spectrum, consistent with the canonical architecture \cite{c09:QG223}.
\end{proof}

\section{Difference Duality (\texorpdfstring{$\rho,\psi$}{rho, psi})}
\label{c09:sec:duality}

The Difference duality structures quantum mechanics: the two faces of Difference --- the
scalar $\rho$ and the tensor $\psi$ --- carry, respectively, the amplitude and the
gravitational/tensor content of the state. The duality \cite{c09:QG286} establishes that
$\rho$ and $\psi$ are the trace and traceless faces of the one Difference
(Chapter~\ref{c01:chap:difference}, Theorem~\ref{c01:thm:duality}); the scalar face is the
counting measure whose share is the amplitude magnitude squared
(Section~\ref{c09:sec:amplitudes}), and the tensor face is the Weyl content read against the
reference $\eta$ (Chapter~\ref{c02:chap:eta}, Theorem~\ref{c02:thm:tensor-read}). The duality
is invariant --- it changes meaning, not numbers \cite{c09:QG301}: the quantum structure is
consistent with it, and no number changes under the duality reinterpretation.

\section{Amplitudes and Probabilities}
\label{c09:sec:amplitudes}

\begin{theorem}[The amplitude magnitude is derived]
\label{c09:thm:amplitude-derived}
The \emph{amplitude magnitude} of a state at generation $k$ is the square root of its
counting-measure share: $|\psi_k| = \sqrt{\rho_k}$, where
$\rho_k = \mu^k/S$ is the actualization share, $\mu$ the branching ratio of the
Galton--Watson process, and $S = \sum_{j<K_{\mathrm{gen}}}\mu^j$, with $K_{\mathrm{gen}}$ the
total number of generations. The amplitude magnitude is derived from
the actualization process:
\begin{equation}
\label{c09:eq:amplitude}
|\psi_k|^2 \;=\; \rho_k \;=\; \frac{\mu^k}{S},
\end{equation}
the counting-measure share of the state. The Born rule $\sum_k|\psi_k|^2 = 1$ holds exactly
by construction, as the normalization of the actualization share: probabilities are the
normalized actualization shares.
\end{theorem}

\begin{proof}
The amplitude-origin result \cite{c09:QG216} derives the amplitude magnitude from the Q-event
process: the path multiplicity to generation $k$ is $\mu^k$, and the normalized share
$\rho_k = \mu^k/S$ is the amplitude magnitude squared. The normalization
$\sum_k|\psi_k|^2 = 1$ is exact by construction for any branching ratio $\mu$; the Born
rule follows from the normalization of the share, so probability is the actualization share
--- the measure, not a separately imposed rule.
\end{proof}

Here $\mu^k$ is the combinatorial path multiplicity of the Galton--Watson tree, namely the
number of root-to-generation-$k$ paths, entering only through the normalized weight
$\rho_k = \mu^k/S$. It is not identified with the conserved Q-event count of
Theorem~\ref{c03:thm:count-conservation}.

\section{Phase Structure}
\label{c09:sec:phase}

\begin{theorem}[The phase is derived from actualization]
\label{c09:thm:phase-derived}
The \emph{quantum phase} of a state at depth $k$ is the circulation phase of the
actualization cycle:
\begin{equation}
\label{c09:eq:phase}
\theta_k \;=\; \frac{2\pi k}{N},
\end{equation}
where $N=96$ is the network periodicity and $k$ the branch depth (the actualization tick).
The phase is fixed by causal ordering and the network periodicity of the $N=96$ attractor.
A quantum state carries exactly two independent real degrees of freedom --- the
magnitude $|\psi| = \sqrt{\rho}$ (the node property) and the phase $\theta$ (the link
property) --- and therefore must be represented as a complex state.
\end{theorem}

\begin{proof}
The phase-origin result \cite{c09:QG220} derives the phase from the Q-event process: causal
ordering fixes the position $k$ (the branch depth, the actualization tick), and the network
periodicity of the circulant ring $C_N$, with $N=96$, fixes the phase quantum
$\Delta\theta = 2\pi/N$ by cycle closure --- $N$ ticks advance $2\pi$, giving uniform
circulation. Hence $\theta_k = 2\pi k/N$ is the circulation phase of the actualization
cycle. The Hilbert-origin result \cite{c09:QG218} establishes the two independent real
degrees of freedom: the magnitude $|\psi| = \sqrt{\rho}$ (from the branching counting
measure, a node property) and the phase $\theta$ (from the $U(1)$ link connection, a link
property). A state with magnitude and phase is a complex state; a real-only state space
would give classical addition with no interference. Hence the complex structure is derived
from the two degrees of freedom.
\end{proof}

The lattice phase $\theta_k = 2\pi k/N$ governs \emph{state evolution}. Mixing and CP
phases are \emph{separate} objects: they are derived from spectral ratios via the
chiral-circulation construction, $\delta = \arcsin(r)$ for a spectral ratio $r$, are
continuous functions of those ratios, and are \emph{not} restricted to the state-phase
lattice.

\section{Interference}
\label{c09:sec:interference}

\begin{theorem}[Interference is phase-dependent]
\label{c09:thm:interference}
Interference in the theory is phase-dependent: the two-path amplitude is
\begin{equation}
\label{c09:eq:interference}
P \;=\; \left|e^{i\theta_1} + e^{i\theta_2}\right|^2
\;=\; 2 + 2\cos(\theta_1 - \theta_2),
\end{equation}
with the phase difference producing the interference pattern. Interference is a derived
consequence of the phase structure: the phase-difference pattern of the
actualization-derived phases.
\end{theorem}

\begin{proof}
The Hilbert-origin result \cite{c09:QG218} establishes the interference formula from the
complex state structure: the two-path amplitude is the sum of the two phase-carrying
amplitudes, whose magnitude squared is $2 + 2\cos(\theta_1 - \theta_2)$. The phase
difference, derived from the actualization cycle (Theorem~\ref{c09:thm:phase-derived}),
produces the interference; a real-only state space would give $P = P_1 + P_2$ with no
interference, so the complex structure is necessary. Hence interference is a derived
consequence of the phase structure, not an independent quantum postulate.
\end{proof}

\section{Measurement as Actualization}
\label{c09:sec:measurement}

\begin{theorem}[Measurement is actualization]
\label{c09:thm:measurement-actualization}
\emph{Measurement} is the collapse of the state to a definite outcome: a Born-weighted
projection identified with the actualization event --- a Q-event is a Born-weighted
projection of the state onto a definite outcome, so no separate collapse postulate or
mechanism is required. Consistent with MONO007 \cite{c09:MONO007}, the measurement result
is presented as the accepted canonical identification: collapse equals actualization. The
theory does not add a new mechanism; it identifies the collapse with the process already
established. No stronger claim is made.
\end{theorem}

\begin{proof}
The complete-quantum-gravity audit \cite{c09:QG223} adjudicates the measurement problem: the
collapse is identified with the Q-event actualization --- a Born-weighted projection to a
definite state. Since the Born weights are the actualization shares
(Theorem~\ref{c09:thm:amplitude-derived}), the projection is weighted by the shares and
performed by the actualization event: measurement is actualization, with no separate
collapse mechanism or postulate.
\end{proof}

\section{Consequences}
\label{c09:sec:qm-consequences}

\begin{lemma}[The quantum dynamics is the actualization flow]
\label{c09:lem:qm-dynamics}
The quantum dynamics of the theory is the actualization flow: the interaction dynamics is
the generator action on the spectral modes, and the Lagrangian is the actualization-flow
action.
\end{lemma}

\begin{proof}
The gauge-dynamics origin \cite{c09:QG243} establishes that the interaction dynamics IS the
generator action on the spectral modes: the D96 gauge generators act on the modes, and an
interaction is the generator's action on the mode. The Lagrangian origin \cite{c09:QG244}
derives the Lagrangian density as the actualization-flow action of the D96 generator
fields, $L = -\frac{1}{4}F^a_{\mu\nu}F^{a\mu\nu} + i\bar{\psi}\gamma^\mu D_\mu\psi -
m\bar{\psi}\psi$. Hence the quantum dynamics is the actualization flow, not an imported
dynamics.
\end{proof}

\begin{theorem}[The conservation structure is universal]
\label{c09:thm:qm-conservation}
The quantum conservation structure is the universal conservation principle: the Born norm,
unitarity, and the Noether currents are manifestations of the one conservation rooted in the
identity of Difference.
\end{theorem}

\begin{proof}
The conservation-principle audit \cite{c09:QG267} establishes that the norm (Born rule),
unitarity, trace, count and Bianchi conservation, and the Noether currents are
manifestations of one deeper principle; count conservation is the definitional identity of
the primitive (Chapter~\ref{c03:chap:actualization},
Theorem~\ref{c03:thm:count-conservation}). Hence the quantum conservation structure is the
universal conservation principle, rooted in the identity of Difference.
\end{proof}

Quantum mechanics is therefore consistent with the canonical architecture: states,
amplitudes, phases, interference, measurement, and dynamics are all derived reads of the
D96 spectrum through the Difference duality and the actualization process
\cite{c09:QG286,c09:QG223}.

\section{Summary}
\label{c09:sec:qm-summary}

This chapter has derived the quantum structure of the theory. Quantum mechanics is
emergent in The Actualization Theory: states are spectral reads of the D96 spectrum
(Theorem~\ref{c09:thm:states-spectral}); the amplitude is the counting-measure share,
$|\psi|^2 = \rho$, with the Born rule exact by construction
(Theorem~\ref{c09:thm:amplitude-derived}); the phase is the circulation phase $2\pi k/N$ of
the actualization cycle, and the state is complex (Theorem~\ref{c09:thm:phase-derived});
interference is the phase-difference pattern (Theorem~\ref{c09:thm:interference});
measurement is the actualization event, with no separate collapse postulate
(Theorem~\ref{c09:thm:measurement-actualization}); and the dynamics is the actualization flow
with universal conservation (Lemma~\ref{c09:lem:qm-dynamics},
Theorem~\ref{c09:thm:qm-conservation}). The duality structures QM
(Section~\ref{c09:sec:duality}): the scalar face carries the amplitude, the tensor face the
gravitational content. The quantum pillar is derived, not postulated. The following
chapters derive gravity and spacetime, the standard model, and cosmology from the same
spectral foundation.

\begin{thebibliography}{99}
\bibitem{c09:QG216} AT-QG Phase 216, \emph{Quantum Amplitude Origin}, Docs/Research.
\bibitem{c09:QG218} AT-QG Phase 218, \emph{Hilbert Space Origin}, Docs/Research.
\bibitem{c09:QG220} AT-QG Phase 220, \emph{Quantum Phase Origin}, Docs/Research.
\bibitem{c09:QG223} AT-QG Phase 223, \emph{Final Quantum Gravity Audit} (Complete QG), Docs/Research.
\bibitem{c09:QG243} AT-QG Phase 243, \emph{Gauge Dynamics Origin}, Docs/Research.
\bibitem{c09:QG244} AT-QG Phase 244, \emph{Lagrangian Origin}, Docs/Research.
\bibitem{c09:QG267} AT-QG Phase 267, \emph{Conservation Principle Audit}, Docs/Research.
\bibitem{c09:QG286} AT-QG Phase 286, \emph{Difference Duality Audit}, Docs/Research.
\bibitem{c09:QG301} AT-QG Phase 301, \emph{Duality Prediction Audit}, Docs/Research.
\bibitem{c09:MONO007} AT-MONO007, \emph{Referee-Readiness Resolution}, Docs/Zenodo/Monograph.
\end{thebibliography}
